LIPAS (Lunar Impact and Predictive Analysis System) is a hybrid forecasting and decision-support platform that combines machine learning with physics-based computational models to predict the Moon's space-weather conditions. These conditions include temperature, radiation, dust activity, micrometeorite impacts, solar storms, and moonquakes. By unifying space-weather data products from NASA and the National Oceanic and Atmospheric Administration (NOAA) alongside OMNI2 data, the LIPAS platform processes the data through a temporal Long Short-Term Memory (LSTM) neural network. Users can view the LIPAS space-weather predictions via an interactive, web-based dashboard for any lunar location.

This study aims to address three research questions. First, it documents how the hybrid model LIPAS was developed and what was learned throughout the process. Second, it studies how operational usability can be improved in real mission scenarios compared with existing prediction systems. Third, this research investigates whether integrating machine learning and physics-based modeling into LIPAS improves prediction accuracy compared with existing single-purpose NASA and NOAA lunar environmental models.
